% ============================================================
% Paper 0 — RBD-Bench unified paper
% Representational-Behavioural Dissociation in LLMs
% Six axes (D1-D6), named routing failures, mechanistic cures
% Merges: paper0_rewrite_v2.tex, F-2 workshop draft, paper0_final.tex
% ============================================================
\pdfoutput=1
\documentclass[11pt]{article}
\usepackage{ACL2023}
\usepackage{times}
\usepackage{latexsym}
\usepackage[T1]{fontenc}
\usepackage[utf8]{inputenc}
\usepackage{microtype}
\usepackage{inconsolata}
\usepackage{amsmath}
\usepackage{float}
\usepackage{graphicx}
\usepackage{multirow}
\usepackage{booktabs}
\usepackage{url}
\usepackage{xcolor}
\usepackage{enumitem}
\usepackage{array}

% Series-consistent macros
\newcommand{\rbd}{\textsc{rbd}}
\newcommand{\tighttable}{\setlength{\tabcolsep}{4pt}\renewcommand{\arraystretch}{1.05}}

% Named mechanisms from F-2
\newcommand{\umg}{Use-Mention Gap}
\newcommand{\Umg}{Use-Mention Gap}
\newcommand{\irb}{Internal Reasoning Buffer}
\newcommand{\Irb}{Internal Reasoning Buffer}
\newcommand{\cc}{Contextual Crystallization}
\newcommand{\Cc}{Contextual Crystallization}
\newcommand{\gesup}{Generative Entrenchment}
\newcommand{\Gesup}{Generative Entrenchment}
\newcommand{\sis}{stance-induced suppression}
\newcommand{\Sis}{Stance-Induced Suppression}
\newcommand{\bet}{binary evaluation trap}
\newcommand{\Bet}{Binary Evaluation Trap}
\newcommand{\ch}{confirmation hysteresis}
\newcommand{\Ch}{Confirmation Hysteresis}
\newcommand{\asa}{API surface anchoring}
\newcommand{\Asa}{API Surface Anchoring}
\newcommand{\ter}{Trajectory Escape Rate}
\newcommand{\TER}{TER}
\definecolor{lightgray}{gray}{0.93}

\title{Representational--Behavioural Dissociation in Large Language Models:\\
Named Routing Failures, Mechanistic Cures, and RBD-Bench}

\author{Elliot Tower}

\begin{document}
\maketitle

% ============================================================
\begin{abstract}
Large language models routinely demonstrate knowledge, rules, and
capabilities in direct probing yet fail to deploy them under the
structural conditions of generation and self-evaluation. We argue that
these failures are not independent bugs but instances of a single
underlying phenomenon: \textbf{representational--behavioural
dissociation} (\rbd{}), adopting the vocabulary of
\citet{vegner2025systematicity}. The model possesses the representational
resources required to succeed but fails to deploy them under the specific
structural conditions of the autoregressive forward pass.

We organize \rbd{} along six axes, each with a named behavioral
failure, a named routing mechanism, and (for D1--D4) a named mechanistic
cure identified via activation steering and head attribution in
small models: (1)~\textbf{D1: Knowledge Activation} --- the
\umg{}, cured by Numeric Recomputation; (2)~\textbf{D2: Compositional
Assembly} --- \irb{} Collapse, cured by Buffer Externalization;
(3)~\textbf{D3: Self-Evaluation} --- \cc{} and Context Lock-Out,
cured by Context Surgery; (4)~\textbf{D4: Trajectory Revision} ---
\Gesup{} and the First Wrong Token Wins failure, cured by Early
Intervention within the pivot-token escape window; (5)~\textbf{D5:
Competence Boundary Awareness} --- confidence does not track accuracy
as constraints tighten; and (6)~\textbf{D6: Structural Verification}
--- binary framing suppresses fine-grained structural analysis even when
the underlying capability is intact.

We introduce \textbf{RBD-Bench}, a benchmark organized around these six
axes with a sub-bench per axis. For D1--D4 we report both large-scale
API behavioral results (11+ models, 4 domains) and small-model
mechanistic localization (head attribution, activation steering,
cross-axis circuit independence). For D5 and D6 we report behavioral
results only; mechanistic localization is planned work. Across 18+
models from 7B to 671B parameters the dissociation is systematic,
model-size-correlated, and largely invariant to standard enhancement
techniques. Five ablation studies document structural aspects of the
failures. The cross-axis steering independence result (within-axis
improves, cross-axis near zero, placebos at baseline) provides the
strongest causal evidence that these are distinct routing circuits.
\end{abstract}

% ============================================================
\section{Introduction}
\label{sec:intro}

Standard benchmarks measure whether a model \textit{can} perform a task
under favorable conditions. They do not measure whether it
\textit{will} under the structural conditions that arise during real
deployment: mid-generation perturbation, multi-step composition,
self-verification, trajectory revision, constraint gradients, and
binary verdict framing.

This paper unifies a class of failures that standard evaluation misses.
In each case the model demonstrably possesses the knowledge or capability
required to succeed---confirmed by a probe condition in which the
capability is elicited---yet systematically fails to deploy it under a
different structural condition. We term this \textbf{representational--%
behavioural dissociation} (\rbd{}), adopting the vocabulary of
\citet{vegner2025systematicity}, who documented cases where end-to-end
neural models exhibit \textit{representational systematicity} (internal
representations satisfy the structural requirements for a task) without
\textit{behavioural systematicity} (actually deploying those
representations across the full space of tasks they license).

The theoretical interpretation draws on two philosophical traditions.
\citet{fodor1983modularity}'s modularity thesis, generalized from
perceptual modules to the generative pathways of autoregressive language
models, holds that parametric knowledge may be informationally encapsulated
from the generative module under specific activation conditions---present
and accessible under explicit metacognitive prompting, yet unavailable
during generation. Quinean holism \citep{quine1951two} accounts for why
committed trajectories resist revision: the generated tokens form a web
of mutual reinforcement that makes any single correction signal
insufficient to overcome the prior commitment.

We organize \rbd{} along six axes, each named by both the behavioral
failure it produces and the mechanism that drives it:

\begin{enumerate}[leftmargin=*, noitemsep, topsep=4pt]
\item \textbf{D1: Knowledge Activation / \Umg{}} (Paper 1): The model
can state a rule when asked but fails to apply it during generation.
Mechanism: binary YES/NO frame routes around the rule-deployment circuit.
Cure: Numeric Recomputation---ask for the answer, not YES/NO.
\item \textbf{D2: Compositional Assembly / \Irb{} Collapse} (Paper 2):
The model knows each sub-step individually but cannot chain them as
depth increases.
Mechanism: the Assembly Gap between rule-description access and
rule-execution capacity. Cure: Buffer Externalization---scaffold
intermediate state into context.
\item \textbf{D3: Self-Evaluation / Context Lock-Out} (Paper 3):
The model generates a correct answer but falsely rejects it when asked
to verify in the same context. Fresh context restores performance.
Mechanism: \Cc{}---prior generation crystallizes evaluative stance.
Cure: Context Surgery---remove committed trace, re-evaluate fresh.
\item \textbf{D4: Trajectory Revision / First Wrong Token Wins}
(Paper 4): Once the model commits to a trajectory, even an explicit
correction signal cannot redirect it.
Mechanism: Web of Commitment---first committed wrong token pulls
subsequent tokens as axiomatic.
Cure: Early Intervention within the pivot-token escape window
($\sim$147 tokens for R1).
\item \textbf{D5: Competence Boundary Awareness} (Paper 5): As
constraints tighten, accuracy drops non-linearly at a cliff point. The
model continues to answer confidently past its own competence boundary.
\item \textbf{D6: Structural Verification / \Bet{}} (Paper 6): Binary
yes/no framing suppresses fine-grained structural analysis, even when
the underlying capability is intact.
\end{enumerate}

For D1--D4 we report both large-scale behavioral sweeps across 11+
models and mechanistic localization results on small models (head
attribution, activation steering, cross-axis circuit independence
controls). D5 and D6 currently have behavioral results only; mechanistic
localization is planned work.

Each axis is documented in a companion paper. This paper provides the
unified account, the benchmark design, the intervention structure,
cross-axis circuit independence evidence, and theoretical discussion.

% ============================================================
\section{RBD-Bench}
\label{sec:bench}

We introduce \textbf{RBD-Bench}, a benchmark organized around the six
axes of \rbd{}. Each axis has its own sub-bench (Knowledge-Activation
Bench, Composition Bench, Context-Lock Bench, Trajectory-Revision Bench,
Competence-Boundary Bench, Structural-Verification Bench) with its own
experimental protocol, datasets, and metrics, but all share a common
structure:

\begin{enumerate}[noitemsep, topsep=2pt]
\item A \textbf{capability probe} that confirms the model possesses the
relevant knowledge or skill.
\item A \textbf{structural condition} under which that capability must
be deployed (mid-generation, multi-step, self-evaluation, etc.).
\item A \textbf{dissociation metric} that quantifies the gap between
probe performance and structural-deployment performance.
\end{enumerate}

\begin{table}[H]
\centering\small\tighttable
\begin{tabular}{p{2.2cm}p{2.0cm}p{1.8cm}p{2.8cm}}
\toprule
\textbf{Axis} & \textbf{Domains} & \textbf{Models} & \textbf{Metric} \\
\midrule
D1: Knowledge Act. & Math, Symbolic rules & 7 models & Recovery rate (\%) \\
D2: Comp. Assembly & Symbolic rules, logic & 6 models & Depth accuracy (\%) \\
D3: Self-Evaluation & SCAN, grammar tasks & 5 models & FRR (\%) \\
D4: Trajectory Rev. & Budget, scheduling, clinical, config & 7 models & TER (\%) \\
D5: Comp. Boundary & Budget, scheduling, clinical, config & 7 models & Confidence--accuracy gap \\
D6: Struct. Verif. & Clinical, code security & 6 models & Recall, F1 \\
\bottomrule
\end{tabular}
\caption{RBD-Bench overview. Each axis has its own sub-bench, domain
coverage, model set, and dissociation metric. D4 and D5 share the
injection-gradient experimental design (A/B/C/D). D6 spans clinical
hallucination detection and code vulnerability detection.}
\label{tab:bench}
\end{table}

The six sub-benches are not separate benchmarks. They are modes of a
single benchmark organized around the same underlying construct. A
model's profile across D1--D6 characterizes the shape of its
dissociation, not just its absolute capability.

Table~\ref{tab:mech-beh} summarizes the mechanism--cure structure for
the four axes with mechanistic localization results.

\begin{table}[H]
\centering\small\tighttable
\begin{tabular}{p{1.5cm}p{2.4cm}p{2.1cm}p{2.1cm}}
\toprule
\textbf{Axis} & \textbf{Behavioral failure} & \textbf{Named mechanism} & \textbf{Named cure} \\
\midrule
D1 & Knowledge not activated & \Umg{} & Numeric Recomputation \\
D2 & Sub-steps don't chain & \Irb{} Collapse / Assembly Gap & Buffer Externalization \\
D3 & Correct output rejected & \Cc{} / Context Lock-Out & Context Surgery \\
D4 & Can't exit bad trajectory & Web of Commitment / First Wrong Token Wins & Early Intervention \\
D5 & Confident past cliff & Competence Boundary failure & (planned) \\
D6 & Binary frame blocks analysis & \Bet{} & Frame Replacement \\
\bottomrule
\end{tabular}
\caption{Behavior, named mechanism, and named cure per axis. D5 and D6
do not yet have mechanistic localization results; only behavioral evidence
is reported for those axes.}
\label{tab:mech-beh}
\end{table}

% ============================================================
\section{D1: Knowledge Activation / \Umg{}}
\label{sec:d1}

\subsection{Behavioral Failure}

The \umg{} is the most tractable failure mode: the model can accurately
\textit{mention} a corrective rule---state it, describe it, identify it
as applicable---yet fails to \textit{use} that same rule during
single-utterance generation. This is an activation deficit, not a
knowledge deficit. The distinction is empirically grounded: when the
model is asked, in a separate metacognitive prompt outside of generation,
whether a given rule applies, it answers correctly. When the model must
apply the rule inside generative flow without being explicitly invited
to state it, it fails.

\subsection{Protocol}

Single-utterance self-correction of synthetically perturbed
chain-of-thought reasoning. A held-out model (LLaMA 3.1 405B) applies
perturbations (decimal shifts, operator swaps, phrase alterations) to
100-token solution stubs; the target model then completes the generation.
Capability probe: the model can identify the perturbation and state the
correct fix in a metacognitive prompt. Dissociation metric: recovery rate
under single-utterance generation.

\subsection{Behavioral Results}

Recovery rates range from 8--16\% for sub-30B models to 87--90\% for
DeepSeek-R1 (671B). The surface signature of gap closure is a class of
discourse markers---\textit{Wait}, \textit{Hold on}, \textit{However},
\textit{Actually}---that signal the model is interrupting its current
reasoning trajectory to initiate a correction. In R1, these pivot tokens
predict recovery in approximately 93\% of cases. In smaller models, the
same markers are produced but carry no predictive value: form without
function.

A secondary finding is \textbf{style-capability coupling}: in QwQ 32B,
the correction pathway activates reliably under on-policy conditions
(completing a stub it generated itself) but fails under off-policy
conditions (completing a stub from a different model), dropping from
66.4\% to 24.0\% recovery on GSM8K. A style-matched ablation confirms
this is not a capability difference but a context gate: when off-policy
stubs are style-matched to QwQ's own generation, recovery returns to
61.2\%.

Domains: math reasoning (GSM8K, GSM-Symbolic, MATH-500) and symbolic
rule induction. The dissociation replicates across both.

\subsection{Mechanistic Localization (Small Models)}

Within-axis activation steering lifts Gemma's chain-verification
discrimination from 0.20 to 0.50 at layer 40. Llama is fully
distributed (0/1024 heads selective). Head attribution: 11/336 Gemma
heads selective (3.3\%), 11/1152 Qwen heads selective (1.0\%).

The mechanism is the binary YES/NO frame: it routes activations around
the rule-deployment circuit, suppressing activation of the correction
pathway. The named cure is \textbf{Numeric Recomputation}---asking for
the answer rather than a YES/NO verdict bypasses the verification circuit
and activates the deployment pathway directly.

% ============================================================
\section{D2: Compositional Assembly / \Irb{} Collapse}
\label{sec:d2}

\subsection{Behavioral Failure}

Models achieve high accuracy on individual sub-steps when tested in
isolation ($>85$\% across model families), yet composed accuracy
collapses sharply as reasoning depth increases: below 60\% at depth 3,
below 35\% at depth 4, below 22\% at depth 6. This is not a knowledge
deficit at the sub-step level. It is an assembly failure: the
sub-competencies are present; they do not compose.

\subsection{Protocol}

Depth-parameterized rule application tasks. Capability probe: the model
applies each rule correctly in isolation. Structural condition: applying
$k$ rules in sequence where each rule's output feeds the next.
Dissociation metric: composed accuracy as a function of depth.

\subsection{Behavioral Results}

The depth wall appears between depths 2 and 4 across model families.
Reasoning-trained models (R1, QwQ) mitigate the wall by generating
extended reasoning traces---explicit buffer externalization: the model
writes intermediate results into the context window, converting an
implicit working-memory problem into an explicit retrieval problem. Trace
length correlates strongly with accuracy at depth for reasoning models
($r > 0.7$) and negligibly for non-reasoning models ($r < 0.2$).

The sharpest evidence: when R1's thinking traces are truncated
mid-assembly, accuracy drops catastrophically from 87\% to 0\% by 50\%
of median trace budget. An account in which traces are epiphenomenal
predicts gradual degradation; we observe a cliff.

Domains: symbolic rule induction (ARI), fictional alphabet ciphers
(FAC), Boolean Dream, Chinese logical reasoning (BLTB). The depth wall
appears at approximately the same composition depth across all four.

\subsection{Mechanistic Localization (Small Models)}

Gemma concentrates 83/336 heads (24.7\%) on composition in layers 1--24.
Falcon3-Mamba is selective in all 64 SSM layers. Mid-layer PCA(5)
logistic regression on Llama layer 16/32 predicts composition failure at
89\% accuracy vs.\ 74\% from output logits alone.

The mechanism is the \textbf{\irb{} Collapse}: the autoregressive forward
pass has limited implicit working memory; multi-step reasoning at depth
exhausts it. The Assembly Gap is the behavioral gap between rule-description
access and rule-execution capacity. The named cure is \textbf{Buffer
Externalization}---scaffolding intermediate state into context converts
the implicit IRB into an explicit retrieval problem.

% ============================================================
\section{D3: Self-Evaluation / Context Lock-Out}
\label{sec:d3}

\subsection{Behavioral Failure}

The Self-Evaluation failure occurs when a model generates a correct
output but falsely rejects it when asked to verify in the same
conversational context. Present the same output in a fresh context (e.g.,
framed as ``a student's answer'') and the model correctly evaluates it.

\subsection{Protocol}

The \textbf{False Rejection Rate} (FRR): the fraction of correctly
generated outputs that the model incorrectly rejects when asked to verify
its own work. Capability probe: third-party evaluation of the same
outputs without generation history. Structural condition: same-context
verification after generation. Dissociation metric: gap between
same-context FRR and fresh-context FRR.

\subsection{Behavioral Results}

Self-evaluation accuracy averages 40--55\%; third-party evaluation yields
75--85\%. The mechanism is \cc{}: the context becomes progressively
crystallized around the model's prior generation, reducing accessible
probability mass for divergent evaluations. The crystallization gradient
is approximately log-linear with trace length, saturating at approximately
500 tokens.

\textbf{Context Surgery}---removing the generation trace from context
before requesting evaluation---restores accuracy to within 5--8
percentage points of third-party evaluation.

At the competence boundary, both same-context and fresh-context
evaluation converge toward constant rejection: the model can no longer
distinguish its correct from incorrect outputs regardless of context.
This absorbing-state regime achieves 0/38 reversals under post-hoc
reframing.

Domains: SCAN grammar translation, symbolic grammar verification (GSV,
Symmball), inductive--abductive dissociation (IAD), second-order rule
shift (KGT). The dissociation replicates across all.

\subsection{Mechanistic Localization (Small Models)}

Qwen has 63/1152 heads (5.5\%) selective for self-vs-neutral framing
with 0\% behavioural FRR---the circuitry distinguishes self from
neutral context but the output ignores this distinction. Sparse Channel
Steering in Falcon3-Mamba: suppressing 14/32 channel groups reduces FRR
from 89\% to 11\% while improving control-task accuracy. This result
holds across 5/5 seeds; placebo controls remain at 81\%~$\pm$~21\%.

The mechanism is \textbf{Context Lock-Out}: prior generation crystallizes
an evaluative stance that locks out divergent evaluation. The named cure
is \textbf{Context Surgery}---removing the committed trace and
re-evaluating in a fresh context.

% ============================================================
\section{D4: Trajectory Revision / First Wrong Token Wins}
\label{sec:d4}

\subsection{Behavioral Failure}

Trajectory Revision is the deepest structural failure mode. Unlike
Knowledge Activation (failure to act on known rules), trajectory revision
failure is the inability to exit a trajectory whose suboptimality the
model has already encoded. \textbf{Latent uncertainty}---entropy over
the top-$k$ token distribution---spikes at the moment of incorrect
commitment, indicating the model has internally signaled that the current
path is suboptimal. Yet the trajectory does not change.

\subsection{Protocol}

The \textbf{Trajectory Escape Rate} (TER): the probability that a model
mid-generation on a suboptimal path escapes to the correct trajectory
when given an explicit one-sentence correction signal. Capability probe:
the model identifies the error when asked directly. Structural condition:
mid-generation on a committed trajectory. Dissociation metric: TER as a
function of generation depth.

The injection-gradient design (A/B/C/D) controls where erroneous
information enters: in the problem premise (\textbf{A}), as a committed
wrong intermediate step (\textbf{B}), as committed steps plus an explicit
contradiction statement (\textbf{C}), or not at all (\textbf{D}).

\subsection{Behavioral Results}

TER decreases monotonically with depth:
\begin{equation}
\mathrm{TER}(d) \propto e^{-\lambda d}
\end{equation}
where $\lambda$ is the entrenchment rate---higher for reasoning models
whose longer traces produce deeper entrenchment. The pivot-token escape
window closes at approximately 147 tokens for R1. After that window,
even reasoning models cannot escape.

Across four domains (manufacturing budget, staff scheduling, clinical
dosing, infrastructure configuration) and seven models, B/C accuracy
collapses to approximately 0--10\% in all but one cell (DeepSeek-V3.2
clinical B~$= 70$\%). Five repair controllers and twelve prompting
interventions all fail. Temperature variation reveals a
detection--correction decoupling: detection rate is modulated by
temperature, but correction rate remains at 0\%.

\subsection{Mechanistic Localization (Small Models)}

Trajectory revision is SFT-resistant. Head attribution: 146/1152 Qwen
heads (12.7\%) selective for trajectory state---the highest specialization
across all four mechanistically localized axes. 44/64 Falcon3-Mamba SSM
layers are selective on trajectory state.

The mechanism is the \textbf{Web of Commitment}: the first committed
wrong token pulls subsequent tokens as axiomatic. This gives rise to the
\textbf{First Wrong Token Wins} failure: once committed, the trajectory
is self-reinforcing. The named cure is \textbf{Early Intervention} within
the pivot-token escape window before the web consolidates.

% ============================================================
\section{D5: Competence Boundary Awareness}
\label{sec:d5}

\textit{Note: D5 currently has behavioral results only. Mechanistic
localization in small models is planned work.}

\subsection{Behavioral Failure}

The injection-gradient design measures where the model's ability actually
cuts off as constraints tighten. Condition A: easy (constraint satisfied,
large margin). Condition D: hard (constraint violated, large margin). The
dissociation is that the model often does not \textit{know} it is near
the cliff. It continues to produce confident answers past its own
competence boundary rather than signaling uncertainty.

\subsection{Protocol}

Capability probe: the model's accuracy on each condition (A/B/C/D)
establishes the actual competence boundary. Structural condition: the
model must produce an answer (forced choice). Dissociation metric: the
gap between the actual competence boundary (where accuracy collapses) and
the model's expressed confidence or willingness to answer.

\subsection{Results}

The A/B/C/D gradient provides an empirical map of the competence boundary
for each model on each task type. Across seven models and four domains,
the confidence--accuracy gap widens at the boundary: the model is most
confidently wrong at the exact point where it is most likely to be wrong.

The prompt-format sweep and abstention results are in
Section~\ref{sec:ablations}. In brief: models with genuine boundary
awareness use the abstention option selectively near the competence
boundary; models without it either never abstain or abstain
indiscriminately. The dissociation metric for D5 is the correlation
between abstention rate and the actual accuracy gradient across A/B/C/D.

Domains: manufacturing budget, staff scheduling, clinical dosing,
infrastructure configuration. Shared with D4 (same injection-gradient
design).

% ============================================================
\section{D6: Structural Verification / \Bet{}}
\label{sec:d6}

\textit{Note: D6 currently has behavioral results only. Mechanistic
localization in small models is planned work.}

\subsection{Behavioral Failure}

Binary yes/no framing suppresses the model's ability to perform
fine-grained structural analysis, even when the underlying capability is
demonstrably intact. This axis spans two structurally independent
domains: clinical hallucination detection and code security vulnerability
detection.

\subsection{Protocol}

Capability probe: the model's performance under an enumeration frame
(``List any claims not supported by the source.'') or diff-format triage
(``Does this commit fix a vulnerability?''). Structural condition: binary
yes/no framing (``Is this hallucinated?'' / ``Is this vulnerable?'').
Dissociation metric: gap between enumeration/diff performance and binary
performance.

\subsection{Results}

\paragraph{Clinical domain.}
Binary prompting achieves recall of 0.324 on real MIMIC-IV hospital
notes. A fact-checker frame achieves recall of 0.838 with precision of
0.812. A student-error frame achieves perfect recall: 86/86 missed
hallucinations recovered. A zero-information nudge---``Please look at the
note and source document more carefully.''---self-corrects 96.7\% of
sampled confident misses.

A $2\times2$ causal experiment across 6~models, 7~datasets, and
2~languages establishes that providing the source document is the only
factor producing consistently large gains (average treatment effect
$= 0.279$, $p < 0.001$), while format choice has negligible average
effect. Even with the source document, performance drops sharply on real
MIMIC-IV hospital notes (F1 falls from 0.851 to 0.484, $-43$\%), and
92.5\% of misses are high-confidence errors (median softmax $= 0.940$).

\paragraph{Code domain.}
Binary prompting achieves F1 of 40--53 across eight benchmarks,
uniformly near chance on balanced splits. Diff-format triage achieves
F1 of 88--100 using the same models. A causal forest analysis across
235~experiments establishes a format average treatment effect of 0.877
(95\%~CI~$[0.767, 0.988]$).

\begin{table}[H]
\centering\small\tighttable
\begin{tabular}{lcc}
\toprule
\textbf{CWE type} & \textbf{Binary F1} & \textbf{Diff F1} \\
\midrule
Buffer overflow (CWE-119/120/121) & 0.44 & 0.91 \\
Use-after-free (CWE-416)          & 0.51 & 0.96 \\
Integer overflow (CWE-190)        & 0.48 & 0.93 \\
Null pointer deref (CWE-476)      & 0.53 & 0.98 \\
Format string (CWE-134)           & 0.41 & 0.89 \\
Other                             & 0.47 & 0.90 \\
\bottomrule
\end{tabular}
\caption{Binary vs.\ diff-format F1 by CWE vulnerability class. The
framing effect is positive across all CWE types. Source:
\citep{tower2026vuln}.}
\label{tab:cwe}
\end{table}

The mechanism of suppression in the code domain is \textbf{\asa{}}:
the binary verification question anchors the model's attention to the
presence of known-dangerous API calls (\texttt{strcpy}: 91.2\% false
positive rate as a spurious confounder; \texttt{memcpy}: 88.4\%) rather
than the structural conditions determining exploitability. Every standard
enhancement technique applied to the binary frame reduces F1 (all 20
tested conditions; mean $\Delta = -0.121$; worst: multi-agent consensus,
$\Delta = -0.930$).

\paragraph{Detection by error type.}

\begin{table}[H]
\centering\small\tighttable
\begin{tabular}{@{}lccc@{}}
\toprule
\textbf{Error type} & $n$ & \textbf{Detection rate} & \textbf{Class} \\
\midrule
Contradicted fact & 21 & 57.1\% & Commission \\
Incorrect fact & 2 & 100\% & Commission \\
\midrule
Location unsupported & 23 & 52.2\% & Omission \\
Word unsupported & 25 & 52.0\% & Omission \\
Time unsupported & 32 & 50.0\% & Omission \\
Condition unsupported & 71 & 43.7\% & Omission \\
Number unsupported & 35 & 40.0\% & Omission \\
Medication unsupported & 48 & 39.6\% & Omission \\
Name unsupported & 19 & 36.8\% & Omission \\
Procedure unsupported & 36 & 27.8\% & Omission \\
\bottomrule
\end{tabular}
\caption{Detection rate by error type under binary framing (MIMIC-IV,
Qwen-2.5-72B). Binary framing handles commission errors and
systematically fails on omission errors. 82.6\% of the 86 confident
false negatives are omission errors. Source: \citep{tower2026detection}.}
\label{tab:types}
\end{table}

Domains: clinical hallucination detection (MIMIC-IV) and code security
vulnerability detection (eight benchmarks).

% ============================================================
\section{Cross-Axis Circuit Independence}
\label{sec:independence}

The four mechanistically localized axes (D1--D4) exhibit circuit-level
independence, providing the strongest causal evidence that these are
distinct routing circuits rather than a single undifferentiated failure.

The cross-axis steering independence experiment tests the specificity of
the steering interventions: does steering within one axis improve that
axis and leave others unchanged? The results are:

\begin{itemize}[noitemsep, topsep=2pt]
\item \textbf{Within-axis steering} improves target-axis performance
consistently (e.g., D3 Context Lock-Out FRR drops from 89\% to 11\%
under Sparse Channel Steering in Falcon3-Mamba).
\item \textbf{Cross-axis steering} produces effects near zero: applying
D3 steering to the D1, D2, or D4 tasks does not change performance on
those tasks.
\item \textbf{Placebo controls} remain at the baseline ($\sim$81\%
$\pm$ 21\%) confirming the within-axis effect is not an artifact of
general activation perturbation.
\end{itemize}

This double dissociation---within-axis improves, cross-axis near
zero---is consistent with four distinct functional circuits and
inconsistent with a unified single mechanism.

The head attribution counts reinforce this picture. Selective head counts
vary substantially across axes (D1: 11/336 Gemma; D2: 83/336 Gemma;
D3: 63/1152 Qwen; D4: 146/1152 Qwen) with negligible overlap in head
identity, suggesting that each routing failure engages a distinct
sub-network.

% ============================================================
\section{Ablation Studies}
\label{sec:ablations}

Beyond the core six-axis protocol we report five ablation studies that
test specific structural claims.

\subsection{Ablation 1: Synthetic MIMIC-IV}

A synthetic clinical note generator matches the word-choice distribution,
sentence length, section structure, and clinical terminology density of
real MIMIC-IV notes but without the underlying clinical reasoning
structure. Models tested on synthetic vs.\ real data under the same D6
protocol.

A causal analysis across 1,390 observations (4 models $\times$ 2
datasets $\times$ 100 items) establishes that the synthetic-vs-real gap
is the strongest single signal:
\begin{itemize}[noitemsep, topsep=2pt]
\item GPT-5.4: 29\% more accurate on synthetic (ATE +0.287, 95\%~CI
[+0.210, +0.364]).
\item GPT-5.4-nothink: 81\% more accurate on synthetic (ATE +0.808,
95\%~CI [+0.746, +0.866]).
\item Qwen-nothink: 69\% more accurate on synthetic (ATE +0.687,
95\%~CI [+0.631, +0.740]).
\end{itemize}
The large ATEs confirm that the failure on real data reflects structural
features in real clinical text that synthetic generation does not capture.

\subsection{Ablation 2: Medical Injection-Gradient}

The injection-gradient design (A/B/C/D) from D4/D5 applied to medical
reasoning tasks. Results replicate the pattern seen in non-medical
domains: models that correctly identify the constraint in isolation
(condition D) collapse to near-zero when a wrong intermediate is
pre-committed (condition B). The thinking--non-thinking divergence is
more pronounced in the medical domain: reasoning-trained models that
produce extended traces are more susceptible to anchoring on the poisoned
intermediate, while instruction-tuned models without extended reasoning
show greater resistance.

\subsection{Ablation 3: Prompt-Format Sweep with Abstention}

A response-format sweep varying what the model produces first (reasoning
trace, verdict, confidence score) and whether the model is given an
explicit abstention option (``I don't know'' / ``I cannot determine'').
Swapped-order conditions also tested.

A pilot ($N=24$) suggested that confidence-elicitation formats produced
a $+33$pp accuracy boost. When expanded to 120 parametric variants
(20 per problem $\times$ 6 problems, spanning the full difficulty
gradient), \textbf{all treatment effects dropped to exactly 0.0\%} with
zero-width confidence intervals. Output format is causally inert for
accuracy and abstention on this domain. The pilot effect was a sampling
artifact.

Models never voluntarily abstain when given the option, regardless of
format. The boundary awareness metric cannot be elicited by format
alone---it requires genuine capability differences between models.

\subsection{Ablation 4: Thinking as Liability}

Controlled ablation: two versions of the \textbf{exact same model}
(Qwen3-235B-Thinking vs.\ Qwen3-235B-Instruct). Same weights, same
architecture, same training data. Only difference: whether the model
generates a chain-of-thought before its answer. $N = 240$ per model
across three domains and four conditions.

Results are stark. On clean conditions (A and D), both models perform
near 100\%. On biased conditions:

\begin{itemize}[noitemsep, topsep=2pt]
\item \textbf{B\_step1} (poisoned intermediate): thinking 1.7\%,
no-thinking 56.7\% ($\Delta = -55.0$pp, $p < 10^{-12}$,
Cohen's $h = -1.45$).
\item \textbf{C\_step2} (model's own wrong output as context): thinking
1.7\%, no-thinking 100\% ($\Delta = -98.3$pp, $p < 10^{-33}$,
Cohen's $h = -2.88$).
\end{itemize}

Overall average treatment effect: $-37.9$pp. On C\_step2, 59 out of 60
paired items are correct only when thinking is disabled. Only 1 item
across all conditions is correct only with thinking enabled.

Most striking: on C\_step2, the thinking model mentions the contradiction
\textbf{100\% of the time} yet still outputs the wrong answer. It
explicitly notes that its prior output was wrong but cannot act on that
awareness. This is the cleanest demonstration of awareness-without-correction
in the series.

\subsection{Ablation 5: Structured Output Heterogeneity}

$N = 1{,}920$ tasks (8 models $\times$ 3 domains $\times$ 10 items
$\times$ 4 conditions $\times$ 2 modes: free text vs.\ JSON-structured).

The effect is highly heterogeneous:
\begin{itemize}[noitemsep, topsep=2pt]
\item \textbf{MiniMax-M27}: JSON improves by +24.2pp (95\%~CI
[+16.1\%, +32.2\%], $p = 3.6 \times 10^{-8}$). On B\_step1, rescues
37.5\% of items from 0\%.
\item \textbf{R1}: JSON improves by +10.0pp (95\%~CI [+4.6\%, +15.4\%],
$p = 4.1 \times 10^{-4}$). On B/C, goes from 0\% to 20\%.
\item \textbf{Qwen-nothink}: JSON \textbf{harms} by $-21.7$pp (95\%~CI
[$-30.1$\%, $-13.2$\%], $p = 1.7 \times 10^{-6}$). On C\_step2, drops
from 100\% to 30\% ($-70$pp).
\item \textbf{GPT-5.4}: unaffected ($\Delta \leq 3.3$pp, n.s.).
\end{itemize}

Structured output is not a universal intervention. It helps models that
benefit from explicit step-by-step computation (MiniMax, R1); it harms
models that already compute correctly without structure (Qwen-nothink).
On C\_step2, JSON mode is catastrophically harmful for Qwen-nothink
($-70$pp)---the structured format forces validation of the poisoned
reasoning chain rather than rejection.

% ============================================================
\section{A Unified Account of Interventions}
\label{sec:unified}

\subsection{Common Failure Structure}

All six failure modes derive from the same architectural constraint: the
autoregressive forward pass is a single, sequential, commitment-prone
computation with no natural mechanism for knowledge activation, working
memory across composition steps, evaluation independence, trajectory
revision, self-monitoring, or fine-grained evidence access under
verdict-first framing. All six are instances of \rbd{}: the model has
the representational resources but the forward pass does not deploy
them. Under \citet{fodor1983modularity}'s modularity framing, the
generated context informationally encapsulates each successive module
from accessing the representation it needs.

\begin{table}[H]
\centering\small\tighttable
\begin{tabular}{p{2.8cm}p{2.5cm}p{3.5cm}}
\toprule
\textbf{Axis} & \textbf{What breaks} & \textbf{What works} \\
\midrule
D1: Knowledge Act. & Knows the rule, doesn't apply it & Numeric Recomputation; activation triggers (RL, pivot elicitation) \\
D2: Comp. Assembly & Rules don't chain at depth & Buffer Externalization (traces, scaffolds, structured state slots) \\
D3: Self-Evaluation & Rejects own correct output & Context Surgery (remove trace before evaluation) \\
D4: Trajectory Rev. & Won't exit bad trajectory & Early Intervention or full restart \\
D5: Comp. Boundary & Confident past its own limit & Abstention option, confidence-first prompts \\
D6: Struct. Verif. & Binary frame blocks analysis & Frame replacement (enumeration, diff-format triage) \\
\bottomrule
\end{tabular}
\caption{Six axes of \rbd{}, what breaks in each case, and the structure
of effective interventions. Each intervention alters the model's
relationship to its own prior context rather than adding new knowledge.}
\label{tab:interventions}
\end{table}

\subsection{The Escalation Trap}

A critical negative result: naive escalation---applying more of the same
kind of engagement to the wrong frame---does not work and typically makes
things worse. The generate-critique-correct pipeline deepens engagement
with the wrong frame and reduces accuracy. Enhancement techniques applied
to the binary frame uniformly reduce F1. Self-correction on entrenched
trajectories deepens commitment to the wrong path.

The reason is consistent: if the source of failure is that the current
context structure blocks access to a capability, then engaging more deeply
with the wrong structure deepens the block. The correct intervention is
structural replacement, not escalation.

\subsection{Effective Intervention Structure}

Effective interventions across all axes share a common property: they
alter the model's relationship to its own prior context.

\begin{enumerate}[noitemsep, topsep=2pt]
\item \textbf{Activating the pathway} (D1): Numeric Recomputation (ask
for the answer directly, not YES/NO); RL modifies the conditions under
which the generative module accesses stored knowledge.
\item \textbf{Externalizing the buffer} (D2): Writing intermediate state
into the context window converts an implicit working-memory problem into
an explicit retrieval problem.
\item \textbf{Removing the crystallizing element} (D3): Context Surgery
removes the generative history before evaluation.
\item \textbf{Intervening early} (D4): Correction must arrive before
entrenchment deepens. After the pivot-token escape window closes
($\sim$147 tokens for R1), only a full restart works.
\item \textbf{Replacing the frame} (D6): Enumeration frames and
diff-format triage activate the structural verification pathway.
\end{enumerate}

\subsection{Structured Verification as the Intervention Class}

The effective interventions across the framing-failure modes share three
structural properties:
\begin{enumerate}[leftmargin=*, noitemsep, topsep=2pt]
\item They request \emph{enumeration} rather than classification
(``list the errors'' vs.\ ``is there an error?'').
\item They presuppose that target conditions exist and ask the model
to characterize them.
\item They activate a sequential claim-level inspection pathway rather
than a holistic verdict pathway.
\end{enumerate}

The student-error frame achieves the clinical recall ceiling (86/86
missed detections recovered). Diff-format triage achieves the code
recall ceiling (F1 88--100). In both cases, the model already had the
evidence; the question type determined whether the structural verification
pathway was activated.

\subsection{Confirmation Hysteresis as Architectural Risk}

\Ch{} is the pipeline-level failure: routing a high-recall detection
frame through a subsequent yes/no confirmation step collapses recall
back to the binary baseline. The student-error frame recovers 100\% of
hallucinations; routing those outputs through a yes/no confirmation
reproduces the original 0\% recall exactly---across both primary
(Qwen-2.5-72B) and holdout (Llama-3.3-70B) models.

This has a specific implication: any pipeline whose terminal verification
step asks a binary question inherits suppression regardless of upstream
sophistication. The fix is not a better binary classifier---it is a
structurally different second stage (claim decomposition, structured
evidence extraction, or any frame that asks ``what specifically is
unsupported?'' rather than ``is this wrong?'').

% ============================================================
\section{Distinguishing \rbd{} from Neighboring Constructs}
\label{sec:distinguish}

\textbf{Hallucination} describes generation of content unsupported by
training data or context \citep{medhallu2025,factscore2023}. \rbd{} is
orthogonal: the failures documented here occur not because the model
generates unsupported content but because it fails to deploy available,
correct capabilities. The high-confidence false negatives in the clinical
domain---92.5\% of misses at median softmax 0.940---are failures of
evaluation, not generation.

\textbf{Miscalibration} describes a mismatch between expressed confidence
and accuracy \citep{guo2017calibration,kadavath2022language}. Several
axes produce miscalibrated outputs, but the cause is structural rather
than a property of the confidence estimator per se. A well-calibrated
model could still exhibit \sis{} if the crystallized context
systematically biases its estimates.

\textbf{Sycophancy} \citep{sharma2023towards,perez2022discovering}
describes models that change answers in response to social pressure.
\rbd{} is structural: the suppression arises from the model's own prior
context, not from user-generated signals.

\textbf{Knowledge gaps} describe failures attributable to absent
knowledge. \rbd{} is definitionally distinguished by the presence of a
probe condition in which the model demonstrates that the relevant
knowledge is present. Every failure mode documented here is accompanied
by such a probe.

\paragraph{The broader measurement problem.}
Binary classification is the near-universal default for LLM evaluation
benchmarks. The results documented here suggest that this protocol
systematically \textit{underestimates} accessible verification capability
on real-world data: it measures commission-detection performance on tasks
whose real-world instances are dominated by structural or omission errors.
LLM-as-judge evaluation that terminates in binary verdicts may be
measuring how well a model responds to a binary question type, not how
well it can perform the underlying verification task.

\paragraph{Atomic decomposition methods.}
FActScore \citep{factscore2023} and SAFE \citep{safe2024} partially
escape the Detection Framing Problem by converting stated claims into
explicit commission checks. This resolves the binary frame problem for
commission errors. However, atomic decomposition does not escape omission
detection: omissions are by definition absent from the claim list that
atomic decomposition is built from. The Detection Framing Problem
therefore persists for the most clinically dangerous error class---missed
information---even under atomic decomposition.

% ============================================================
\section{Discussion}
\label{sec:discussion}

\subsection{Routing Failures as a Unifying Account}

The use--mention gap, context lock-out, and first-wrong-token-wins
failures are three faces of the same phenomenon: the autoregressive
forward pass commits to a generative pathway before it has the
information needed to evaluate that pathway. In the \umg{}, commitment
to a YES/NO answer-frame routes around the rule-deployment circuit.
In context lock-out, commitment to a prior generation crystallizes
the evaluative stance. In trajectory revision, commitment to a wrong
token makes every subsequent token a defense of that token.

Under \citet{fodor1983modularity}'s framing, each of these is a form
of informational encapsulation generalized from perceptual modules to
the generative pathways of autoregressive models: the relevant
representational resource is present but inaccessible from the current
activation context.

\subsection{Latent Uncertainty Present but Unused}

A consistent finding across D1 and D4 is that the model has the right
signal---pivot tokens, latent uncertainty spikes at the moment of
incorrect commitment---but the behavioral system does not act on it. In
D4, latent uncertainty (top-$k$ entropy) spikes at the pivot point yet
trajectory does not change. In D1, pivot tokens predict recovery at 93\%
in R1 but carry zero predictive value in smaller models. The signal is
present in both cases; what differs is whether the behavioral pathway
can act on it.

\subsection{Architecture Determines Cure Type}

The mechanistic localization results reveal that architecture determines
the appropriate cure class. Transformer-based models (Gemma, Qwen,
Llama) show axis-specific head attribution patterns amenable to
attention-head steering. SSM-based models (Falcon3-Mamba) show
axis-specific SSM channel selectivity amenable to sparse channel steering.
The cure family (attention steering vs.\ channel steering) is
architecture-determined, but the axis specificity---and thus the
cross-axis independence of the cure---holds across both architecture
classes.

\subsection{The Autoregressive Bottleneck}

All six failure modes trace to a single architectural feature: the
autoregressive forward pass is a commitment machine. It generates tokens
conditioned on all prior tokens, with the KV-cache structure making prior
tokens exert a persistent, non-revisable attentional influence. This
commitment structure is the source of coherent generation; it is also the
structural bottleneck that makes each of the six \rbd{} failure modes
possible. Effective interventions do not work around this bottleneck by
adding compute (escalation fails) but by changing the structure of the
computation (frame replacement, buffer externalization, context surgery,
early intervention).

\subsection{Fodor's Modularity as Interpretive Lens}

We treat Fodor's modularity thesis as an interpretive lens, not a
mechanistic claim. The thesis holds that perceptual modules are
informationally encapsulated: they process their inputs without access
to background beliefs even when those beliefs would change the output.
We generalize this to the generative pathways of autoregressive models:
parametric knowledge may be encapsulated from the generative module under
specific activation conditions. The encapsulation is not permanent; it
is a function of the activation context. Explicit metacognitive prompting
can break it. RL against verifiable outcomes appears to relax it during
generation \citep{deepseek2025incentivizing,kimi2025scaling}. This
generalization is behavioral; we make no mechanistic claim about attention
weights or hidden states beyond what the intervention evidence supports.

\subsection{Theoretical Connections}

The Self-Evaluation failure (D3) has a self-referential structure: the
model is asked to evaluate its own output from within the same context
that produced it. The crystallization effect means the context cannot
neutrally assess what it created. The entrenchment effect (D4) has a
structural parallel to Quinean confirmation holism
\citep{quine1951two}: the generated tokens form a web of mutual
reinforcement that makes any single correction signal insufficient to
overcome prior commitment. Both are interpretive parallels, not identity
claims.

% ============================================================
\section{Related Work}
\label{sec:related}

The representational--behavioural distinction was surveyed by
\citet{vegner2025systematicity}, who documented cases of representational
systematicity without behavioural systematicity in end-to-end
architectures. Our contribution is to show that this dissociation is
systematic across six axes, to name the mechanisms and cures, and to
localize the circuits mechanistically for four of the six axes.

Chain-of-thought prompting \citep{wei2022chain} and scratchpad methods
\citep{nye2022show} showed that making intermediate computation explicit
improves multi-step reasoning. Our D2 results confirm this and provide a
sharper account: the improvement comes from buffer externalization, not
from better rule reminders. The scaffold ablation (rule scaffolds alone:
0\% accuracy at depth 3; full scaffolds: recovery to within 8\% of
baseline) makes the distinction sharp.

Self-correction literature \citep{madaan2023selfrefine,shinn2023reflexion,
huang2024large} has documented that models often fail to correct their
own outputs. Our D3 results provide a structural explanation: same-context
verification is contaminated by the generative history. Our D4 results
show that the problem is deeper than retrospective self-correction---it
operates during generation itself.

Reinforcement learning against verifiable outcomes
\citep{deepseek2025incentivizing,kimi2025scaling} has shown that simple
RL elicits self-evaluating behavior. Our D1 results confirm this and show
that the correction pathway is latent in base models but reliably
inaccessible without RL or explicit activation triggers.

% ============================================================
\section{Conclusion}
\label{sec:conclusion}

Six structural failure modes in autoregressive language models have been
characterized and unified under a single account: \rbd{}. For four axes
(D1--D4) we name both the behavioral failure and the routing mechanism,
identify an effective cure, and provide mechanistic localization in small
models via head attribution and activation steering. For two axes (D5,
D6) we document behavioral failures with replicated results across 6+
models and two independent domains; mechanistic localization is planned
work.

The cross-axis circuit independence result---within-axis steering
improves, cross-axis near zero, placebos at baseline---provides the
strongest causal evidence that these are distinct routing failures rather
than a single undifferentiated failure mode.

The practical implication is that efforts to improve model performance on
reasoning, verification, and detection tasks should treat the structure of
task framing as a first-class design parameter. Naive escalation is not
only ineffective but systematically harmful. The correct intervention in
each case is structural: activate the pathway, externalize the buffer,
remove the crystallizing element, intervene early, or replace the
verdict-first frame with a structured search frame.

\section*{Limitations}

All D1--D4 mechanistic claims are localized in small models (Gemma,
Qwen, Llama, Falcon3-Mamba) and may not generalize to larger API-only
models. D5 and D6 currently have behavioral results only; mechanistic
localization is planned work. All failure mode accounts are behavioral
in origin. The \irb{} is a behavioral hypothesis without direct
mechanistic confirmation. The trench model is an interpretive
abstraction. Latent uncertainty is operationalized as top-$k$ entropy,
an indirect proxy. The pivot-token escape window ($\sim$147 tokens for
R1) is empirically identified and has not been theoretically derived.
The crystallization gradient saturation point ($\sim$500 tokens) is
identified empirically. The clinical MIMIC-IV analysis is based on 200
annotated notes, sufficient for the framing analysis but underpowered
for subgroup analysis.

\section*{Ethics Statement}

The failure modes documented here---particularly D3 (Self-Evaluation),
D4 (Trajectory Revision), and D6 (Structural Verification)---have direct
implications for clinical AI and security systems deployed at scale.
Binary verification pipelines that report high confidence on outputs they
have not verified at the claim level do not possess the accuracy they
imply. The use of the same evaluative frame in first and second stages of
a pipeline is not a conservative default; it is a systematic risk. We
encourage practitioners to treat task framing as a safety-critical
parameter and to evaluate on real-world data, not only controlled
benchmarks.

\bibliography{geb_series_unified}
\bibliographystyle{acl_natbib}

%================================================================
% APPENDIX
%================================================================
\appendix

\section{Large-Scale API Model Sweep}
\label{app:api-sweep}

\subsection{Trajectory-Revision Bench}
\label{app:trb}

Full A/B/C/D results across 11 models $\times$ 3 domains $\times$ 4
conditions $= 10{,}560$ items. Table~\ref{tab:trb-archetypes} documents
four failure archetypes identified across models.

\begin{table}[H]
\centering\small\tighttable
\begin{tabular}{p{3.2cm}p{5.0cm}}
\toprule
\textbf{Archetype} & \textbf{Description} \\
\midrule
Internal over-committer & Generates long traces that entrench wrong answers (R1 prototype) \\
Silent follower & Accepts injected errors without flagging them (GPT-5.4 prototype) \\
Surface detector & Detects but does not correct committed errors (GPT-4.1 prototype) \\
Injection-immune & Resists injection on clean items, high TER (Nemotron prototype) \\
\bottomrule
\end{tabular}
\caption{Four failure archetypes in Trajectory-Revision Bench.}
\label{tab:trb-archetypes}
\end{table}

\subsection{Context-Lock Bench}
\label{app:clb}

SL-SCAN FRR results, Symmball FRR under 3 conditions, gold-standard
verification near-miss rejection, and GPT-5.4 think vs.\ no-think
absorbing vs.\ graded analysis.

\subsection{Composition Bench}
\label{app:cob}

Depth accuracy by model (full table) and scaffold experiment results.

\subsection{Knowledge-Activation Bench}
\label{app:kab}

Perturbation recovery across 7 models $\times$ 4 math benchmarks
(GSM8K, GSM-Symbolic, MATH-500, ARI).

% ============================================================
\section{Small-Model Mechanistic Analysis}
\label{app:mechanistic}

\subsection{Full Head Attribution Tables}
\label{app:heads}

\begin{table}[H]
\centering\small\tighttable
\begin{tabular}{lcccc}
\toprule
\textbf{Model} & \textbf{D1 heads} & \textbf{D2 heads} & \textbf{D3 heads} & \textbf{D4 heads} \\
\midrule
Gemma & 11/336 (3.3\%) & 83/336 (24.7\%) & --- & --- \\
Qwen  & 11/1152 (1.0\%) & --- & 63/1152 (5.5\%) & 146/1152 (12.7\%) \\
Llama & 0/1024 (0\%, fully distributed) & --- & --- & --- \\
Falcon3-Mamba (SSM) & --- & 64/64 layers & 14/32 ch.\ groups & 44/64 layers \\
\bottomrule
\end{tabular}
\caption{Selective heads (or SSM layers/channel groups) by model and
axis. ``---'' denotes axis not tested on this model in the mechanistic
study. Llama's 0/1024 result for D1 indicates full distribution, not
absence of effect.}
\label{tab:head-attr}
\end{table}

\subsection{Steering Method Details}
\label{app:steering}

Four methods were tested:
\begin{enumerate}[noitemsep, topsep=2pt]
\item \textbf{Global mean diff}: steering direction = mean activation
(correct) $-$ mean activation (incorrect) across all heads and layers.
\item \textbf{Within-rule mean diff}: direction computed within matched
rule pairs only.
\item \textbf{Contrastive direction}: principal component of the
correct/incorrect contrast in activation space.
\item \textbf{Interpolation}: linear interpolation between the
model's activation and the target activation state.
\end{enumerate}

Cross-axis independence was tested using all four methods; the
within-axis-only improvement result holds for all four.

\subsection{Wilson Confidence Intervals}
\label{app:wilson}

All proportion estimates use Wilson confidence intervals
\citep{wilson1927probable}. The placebo control result
(81\%~$\pm$~21\%) uses the Wilson interval for the proportion of
unchanged items under placebo steering.

\subsection{Seed Robustness and Placebo Controls}
\label{app:seeds}

D3 Context Surgery steering result (FRR 89\% to 11\%) replicates
across 5/5 independent seeds. Placebo controls---applying a random
steering direction of the same magnitude as the within-axis
direction---produce FRR of 81\%~$\pm$~21\%, confirming that the
improvement is specific to the axis-appropriate direction.

% ============================================================
\section{Full Methodology}
\label{app:methods}

\subsection{Context Surgery Procedures}
\label{app:surgery}

Context Surgery for D3 removes the generation trace from the KV-cache
before requesting evaluation. In practice this is implemented as a fresh
API call with the generation trace omitted from the context. Three
variants were tested: (1) remove entire assistant turn, (2) remove
assistant turn but retain the final output as a quoted ``student answer'',
(3) paraphrase the assistant turn before requesting evaluation. All three
achieve within 5--8 percentage points of third-party evaluation.

\subsection{TER Measurement Protocol}
\label{app:ter}

$\mathrm{TER}(d) = P(\text{correct trajectory} \mid \text{correction
signal}, \text{depth } d)$. Exponential decay fit: $\mathrm{TER}(d)
\propto e^{-\lambda d}$ where $\lambda$ is estimated per model via
maximum likelihood on the empirical TER curve. The pivot-token escape
window (depth at which the fitted TER curve crosses 0.05) is 147 tokens
for R1 at the median.

\subsection{Injection-Gradient Templates}
\label{app:injection}

\begin{itemize}[noitemsep, topsep=2pt]
\item \textbf{Condition A}: Problem premise contains erroneous assumption
with large violation margin (easy to detect if constraint is checked).
\item \textbf{Condition B}: Problem is clean; a committed wrong
intermediate step is pre-filled in the reasoning trace.
\item \textbf{Condition C}: Condition B plus an explicit one-sentence
contradiction statement appended after the committed step.
\item \textbf{Condition D}: Control; no injection; problem and trace
are correct.
\end{itemize}

\subsection{Fine-Tuning Details}
\label{app:finetune}

QLoRA fine-tuning on Llama-3.1-8B for the SFT-resistance ablation in D4:
rank 16, alpha 32, 4-bit quantization, 3 epochs on the trajectory
revision training set ($N=1200$ items, balanced across A/B/C/D).
Post-SFT TER decline rate $\lambda$ is not significantly reduced
relative to the zero-shot baseline ($p = 0.41$, Wilcoxon signed-rank
test), confirming the SFT-resistance claim.

% ============================================================
\section{Per-Axis Full Results}
\label{app:full-results}

Headline tables with pointers to companion papers for each axis.

\begin{itemize}[noitemsep, topsep=2pt]
\item D1: Recovery rates by model and dataset; full pivot-token
predictive-value table. See \citet{tower2026usemention}.
\item D2: Depth accuracy by model, domain, and depth; full scaffold
ablation. See \citet{tower2026assembly}.
\item D3: FRR by model, domain, and context condition; full Context
Surgery results. See \citet{tower2026lockout}.
\item D4: TER by model, domain, depth, and condition; full repair
controller results. See \citet{tower2026sunkcost}.
\item D5: Confidence--accuracy gap by model and domain; full
prompt-format sweep. Companion paper in preparation.
\item D6: Binary vs.\ alternative-frame performance by model and
benchmark. See \citet{tower2026detection,tower2026vuln}.
\end{itemize}

% ============================================================
\section{Experimental Setup}
\label{app:setup}

\begin{table}[H]
\centering\small\tighttable
\begin{tabular}{p{3.5cm}p{1.5cm}p{1.5cm}p{2.5cm}}
\toprule
\textbf{Model} & \textbf{Params} & \textbf{Provider} & \textbf{Quantization} \\
\midrule
DeepSeek-R1 & 671B & DeepSeek & FP8 \\
DeepSeek-V3.2 & 671B & DeepSeek & FP8 \\
GPT-5.4 & --- & OpenAI & Cloud API \\
GPT-4.1 & --- & OpenAI & Cloud API \\
Qwen3-235B-Thinking & 235B & Alibaba & BF16 \\
Qwen3-235B-Instruct & 235B & Alibaba & BF16 \\
Qwen-2.5-72B & 72B & Alibaba & BF16 \\
Llama-3.3-70B & 70B & Meta & BF16 \\
Llama-3.1-8B & 8B & Meta & BF16 \\
MiniMax-M27 & 27B & MiniMax & BF16 \\
Gemma-2-27B & 27B & Google & BF16 \\
QwQ-32B & 32B & Alibaba & BF16 \\
Nemotron & --- & NVIDIA & Cloud API \\
Command-R7B & 7B & Cohere & BF16 \\
Nemo-12B & 12B & Mistral & BF16 \\
Falcon3-Mamba & 7B & TII & BF16 \\
\bottomrule
\end{tabular}
\caption{Complete model list with parameters, providers, and
quantization used in the large-scale sweep. ``---'' denotes
proprietary models where parameter count is not publicly disclosed.}
\label{tab:models}
\end{table}

\section{Synthetic Diff Construction for Function-Level Inputs}
\label{app:synth-diff}

The diff-format triage prompt requires a commit diff. When only a
function is available (no commit history), a synthetic diff is
constructed: present the function as the ``before'' state and a no-op
identical function as the ``after'' state, then ask whether the diff
reveals that a vulnerability was present in the before version.

\noindent\textbf{Synthetic diff prompt.} \textit{The following is a
synthetic commit diff. The ``before'' version is the function under
analysis; the ``after'' version is an identical no-op copy. Does this
diff reveal that a vulnerability was present in the before version? If
yes, describe the vulnerability type and the structural conditions that
make it exploitable.}

This construction preserves the diff-format activation signal---the
model reasons about a ``fix'' that was never applied---while remaining
applicable to codebases without commit history. CWE formal definition
injection and knowledge-graph confounder injection apply identically to
synthetic diffs. Source: \citep{tower2026vuln}.

\end{document}
